\section{Preliminaries}\label{sec:prelim}

Throughout, $\A=\{0,1\}^2$, and a cell $(x,y)\in\A$ is also written as the integer $2x+y\in\{0,1,2,3\}$; thus $3=(1,1)$ and $0=(0,0)$. The left shift is $\sh(c)_i=c_{i+1}$. A \emph{block map of radius $R$} is a map $g\colon\A^{\Z}\to\A^{\Z}$ of the form $g(c)_i=\gamma(c_{i-R},\dots,c_{i+R})$ for a local rule $\gamma\colon\A^{2R+1}\to\A$; we identify $g$ with the lookup table of $\gamma$, which has $4^{2R+1}$ entries of two bits each. A block map of radius $R$ is also a block map of every larger radius, so ``radius at most $R$'' and ``radius $R$'' describe the same set of maps. By Curtis--Hedlund--Lyndon, $\End(\A^{\Z})$ is the monoid of block maps and $\Aut(\A^{\Z})$ the group of invertible ones. We write
\[
\Cent_R(F)=\{g\ \text{a block map of radius}\le R:\ g\circ F=F\circ g\},\qquad \Cent(F)=\bigcup_R\Cent_R(F),
\]
for the centralizer of $F$ in $\End(\A^{\Z})$, and $\Cent^{\times}_R(F)$, $\Cent^{\times}(F)$ for the invertible elements.

For an elementary rule $f$, identified with its Wolfram number $\sum_{a,b,c}f(a,b,c)2^{4a+2b+c}$, the automaton $F_f$ is defined by \eqref{eq:F}. It has radius $1$, and $F_f^b$ has radius at most $|b|$, so $\sh^aF_f^b\in\Cent_{|a|+|b|}(F_f)$.

\begin{definition}
The \emph{spacetime subshift} of a reversible cellular automaton $F$ on $\A^{\Z}$ is
\[
X_F=\bigl\{(F^n(c))_{n\in\Z}:\ c\in\A^{\Z}\bigr\}\subset\A^{\Z^2},
\]
with the $\Z^2$-action generated by the horizontal shift $\sh$ and the vertical shift $\tau$, $(\tau z)_{n}=z_{n+1}$. An \emph{automorphism} of $X_F$ is a homeomorphism of $X_F$ commuting with both.
\end{definition}

\begin{lemma}\label{lem:spacetime}
For reversible $F$ the map $c\mapsto(F^n(c))_{n\in\Z}$ is a homeomorphism $\A^{\Z}\to X_F$ conjugating $(\sh,F)$ to $(\sh,\tau)$, and it induces an isomorphism
\[
\Aut(X_F)\;\cong\;\Cent^{\times}(F)=\{g\in\Aut(\A^{\Z}):\ gF=Fg\}.
\]
Under this isomorphism $\langle\sh,\tau\rangle$ corresponds to $\langle\sh,F\rangle$.
\end{lemma}

\begin{proof}
The map is continuous, injective (the $0$-th row recovers $c$), and surjective onto $X_F$ by definition, hence a homeomorphism of compact spaces; it carries $\sh$ to $\sh$ and $F$ to $\tau$ since $F^{n}(F(c))=F^{n+1}(c)$. A homeomorphism of $\A^{\Z}$ commuting with $\sh$ is a block map by Curtis--Hedlund--Lyndon, and commuting with $\tau$ means commuting with $F$.
\end{proof}

Thus determining $\Aut(X_F)$ means determining which invertible block maps commute with $F$, and $\Cent(F)$ is the monoid of endomorphisms of the spacetime subshift. The non-invertible elements of $\Cent(F)$ matter too: a block map that commutes with $F$ but is not injective is a non-surjective endomorphism of $X_F$, and the constant map onto a fixed point of $F$ is always one. For every rule with $f(0,0,0)=0$ and $f(1,1,1)=0$, which includes $30$, $54$, $90$ and $110$, both $(0,0)^{\Z}$ and $(1,1)^{\Z}$ are fixed by $F_f$, so $\Cent_R(F_f)$ always contains these two constants; for rule $150$, which has $f(1,1,1)=1$, only $(0,0)^{\Z}$ is fixed.


\subsection{Symmetries}
Three involutions of $\A^{\Z}$ act on second-order automata: the cellwise complement $\kappa(x,y)=(1-x,1-y)$, the reflection $\mathrm{rev}(c)_i=c_{-i}$, and the time reversal $\rho(x,y)=(y,x)$, applied in every cell. The first two are the Boolean conjugation and the reflection of \cite[Section~2]{Takesue1989}, and that exchanging the two rows reverses time is noted in \cite[p.~2500]{Takesue1987}.

\begin{lemma}\label{lem:classes}
For every elementary rule $f$, $\kappa F_f\kappa=F_{\kappa f}$ with $\kappa f(a,b,c)=f(1-a,1-b,1-c)$; $\mathrm{rev}\,F_f\,\mathrm{rev}=F_{f^r}$ with $f^r(a,b,c)=f(c,b,a)$, and $\mathrm{rev}\,\sh\,\mathrm{rev}=\sh^{-1}$; and $\rho F_f\rho=F_f^{-1}$. For $\theta\in\{\kappa,\mathrm{rev}\}$ the map $g\mapsto\theta g\theta^{-1}$ preserves the radius of block maps, maps $\Cent_R(F_f)$ onto $\Cent_R(F_{f'})$ and $\Aut(X_{F_f})$ onto $\Aut(X_{F_{f'}})$, where $f'=\kappa f$ or $f^r$, and maps $\langle\sh,F_f\rangle$ onto $\langle\sh,F_{f'}\rangle$.
\end{lemma}

\begin{proof}
By \eqref{eq:F}, $F_f(\kappa c)_i=(1-y_i,\,1-x_i\oplus f(1-y_{i-1},1-y_i,1-y_{i+1}))$, and applying $\kappa$ gives $(y_i,\,x_i\oplus\kappa f(y_{i-1},y_i,y_{i+1}))$. Similarly $\mathrm{rev}\,F_f\,\mathrm{rev}(c)_i=F_f(\mathrm{rev}\,c)_{-i}=(y_i,\,x_i\oplus f(y_{i+1},y_i,y_{i-1}))$. And $F_f(\rho c)_i=(x_i,\,y_i\oplus f(x_{i-1},x_i,x_{i+1}))$, whose image under $\rho$ is $F_f^{-1}(c)_i$. Conjugation by a cellwise map or by $\mathrm{rev}$ preserves the radius, and the remaining statements follow, since $\kappa$ commutes with $\sh$ and $\mathrm{rev}$ conjugates $\sh$ to $\sh^{-1}$.
\end{proof}

So $\Aut(X_{F_f})$, and every statement below about block maps of bounded radius, depends only on the class of $f$ under $\kappa$ and reflection. There are $88$ classes \cite{Takesue1987}, \cite[Table~I]{Takesue1989}.

\subsection{Diagrams}
For $c\in\A^{\Z}$ the \emph{diagram} of $c$ is $Y_c\colon\Z^2\to\{0,1\}$, where $Y_c(t,i)$ is the $y$-bit of $F^t(c)_i$. The $x$-row of $F(d)$ is the $y$-row of $d$, so the $x$-bit of $F^t(c)_i$ is $Y_c(t-1,i)$, and $c_i=(Y_c(-1,i),Y_c(0,i))$. By \eqref{eq:F} every diagram satisfies the \emph{spacetime relation}
\begin{equation}\label{eq:rel}
Y(t+1,i)\oplus Y(t-1,i)=f\bigl(Y(t,i-1),Y(t,i),Y(t,i+1)\bigr)\qquad\text{for all }(t,i)\in\Z^2,
\end{equation}
so two consecutive rows determine all rows, and $c\mapsto Y_c$ is a bijection from $\A^{\Z}$ onto the solutions of \eqref{eq:rel}. For $t\ge0$, $Y_c(t,i)$ depends only on $c_{[i-t,i+t]}$. Moreover $Y_{\sh c}(t,i)=Y_c(t,i+1)$, $Y_{F(c)}(t,i)=Y_c(t+1,i)$ and $Y_{\rho c}(t,i)=Y_c(-1-t,i)$, the last because $F^t(\rho c)=\rho F^{-t}(c)$.

\subsection{Blocking words, window graphs and eigenfields}\label{sec:blocking}
Let $F$ be a cellular automaton of radius $1$ on $\A^{\Z}$. A word $w$ on an interval $I$ of cells is \emph{blocking} for $F$ (cf.\ the proof of \cite[Theorem~3]{Kurka1997}; the term is from \cite{BT2000}) if there is a cell $i_0\in I$ such that $F^n(c)_{i_0}=F^n(c')_{i_0}$ for all $n\ge0$ and all configurations $c,c'$ that carry $w$ on $I$. A word that is fixed by $F$ in every context is blocking.

Let $F$ be a bijective cellular automaton on $\A^{\Z}$. For $(k,p)\in\Z^2$ with $p\ge1$ and $M\ge0$, let $W^{(k,p)}_M$ be the graph on $\A^{M+1}$ with an edge from each word to each of its possible images under $\sh^kF^p$, the cells outside the window ranging over all values: there is an edge from $c_{[0,M]}$ to $\sh^kF^p(c)_{[0,M]}$ for every $c\in\A^{\Z}$. The graph $W^{(1,1)}_M$ is the light-speed window graph $\Lambda_M$ used for the permutive rules \cite{PaperB}, and $W^{(0,1)}_M$ is a zero-speed graph. We write $(\Lambda^{\pm})$ for the condition that $W^{(1,1)}_M$ and $W^{(-1,1)}_M$ are strongly connected for every $M$, and we call the condition that $W^{(0,1)}_M$ is strongly connected for every $M$ the \emph{zero-speed condition}. A nonconstant block map $\vartheta\colon\A^{\Z}\to B^{\Z}$ onto sequences over a finite alphabet $B$ with $\vartheta\circ F^p=\sh^{-k}\circ\vartheta$ is an \emph{eigenfield of speed $k/p$}, and a \emph{frozen field} if $k=0$; composing with the indicator of a letter shows that one exists if and only if one exists with $B=\{0,1\}$.

\begin{lemma}[{\cite[Lemma~2.5]{PaperE}}]\label{cl:lem:eigen}
Let $F$ be a bijective cellular automaton on $\A^{\Z}$ and $(k,p)\in\Z^2$ with $p\ge1$. The following are equivalent.
\begin{enumerate}
\item[(1)] The graph $W^{(k,p)}_M$ is strongly connected for every $M\ge0$.
\item[(2)] There is no nonconstant block map $\vartheta\colon\A^{\Z}\to\{0,1\}^{\Z}$ with $\vartheta\circ F^p=\sh^{-k}\circ\vartheta$.
\item[(3)] Every continuous function $v\colon\A^{\Z}\to\{0,1\}$ with $v\circ\sh^kF^p=v$ is constant.
\end{enumerate}
\end{lemma}

By the lemma, $(\Lambda^{\pm})$ says that $F$ has no eigenfield of speed $\pm1$, and the zero-speed condition says that $F$ has no frozen field. A frozen field is a local conservation law in the sense of Takesue \cite{Takesue1987}, \cite[Section~4]{Takesue1989}, and $W^{(0,1)}_M$ is the graph of the transition matrix on words of $M+1$ cells considered there; the equivalence of (1) and (2) for $(k,p)=(0,1)$ is stated in \cite[Section~4]{Takesue1989} for the second-order elementary automata (see also \cite[Proposition~4.1]{HT1991}).
